\chapter{Espectro global del paso discreto}
\label{ch:brillouin-global}

\section{Estado al cierre de Fase 0}

Quedo bien establecido que la cuasienergia del paso discreto es circular y
que la familia de esquinas exhibe ya estructura no trivial; quedo claro que
en tiempo discreto los sectores $0$ y $\pi$ pueden coexistir. No quedo
demostrada todavia una clasificacion exhaustiva de \emph{todos} los puntos
especiales del espectro global.

\section{Tareas pendientes}

\begin{enumerate}[label=(B\arabic*),leftmargin=*]
  \item Catalogar puntos especiales del espectro de \Uop\ sobre la zona de
        Brillouin del paso ordenado minimo.
  \item Decidir cuales son fijos bajo el grupo discreto exacto del paso, y
        cuales son artefactos del esquema reciproco.
  \item Producir un test \texttt{validators/brillouin\_global.py} (sympy +
        z3) que confirme la clasificacion para tamanos pequenos.
\end{enumerate}
